\documentclass[a4paper,12pt]{ltjsarticle}
\usepackage[top=25mm,bottom=25mm,left=25mm,right=25mm]{geometry}
\usepackage{graphicx}
\usepackage{float}
\usepackage{booktabs}
\usepackage{array}
\usepackage{enumitem}
\usepackage{hyperref}
\usepackage{caption}

\begin{document}

\begin{center}
  {\Large \textbf{第8回議事録}}\\[8pt]
  {\large チーム4}
\end{center}

\vspace{4mm}

%-------------------------------------------------------------------
\section*{1.基本情報}
\begin{tabular}{ll}
  \toprule
  日時     & 2026年6月10日（水） \\
  会議場所 & 731教室 \\
  目的     & 進捗報告 \\
  チーム番号 & 4 \\
  議事録作成者 & 細澤悠真 \\
  \bottomrule
\end{tabular}

\vspace{4mm}
\noindent\textbf{参加者}

\begin{tabular}{ll}
  2年 & 山口汐里，佐藤夏美，成毛海人 \\
  3年 & 細澤悠真，武本龍，酒井睦基 \\
\end{tabular}

%-------------------------------------------------------------------
\section*{2.議題一覧}
\begin{enumerate}[label=議題\arabic*：]
  \item 議論して決まった内容
  \item 計画からの変更点
  \item 個人毎の次回までに実施する内容
\end{enumerate}

%-------------------------------------------------------------------
\section*{3.議題ごとの内容}

\subsection*{議題①：議論して決まった内容}
\subsubsection*{決定事項}
\begin{itemize}[leftmargin=1.5em]
  \item フルート・クラリネット・オルガン・ドラムの4楽器による輪唱が動作することを確認した（佐藤が当日の動作動画を撮影・共有済み）。実行ファイルは山口がGitHubの\texttt{main}へアップロードした最新の\texttt{楽器.pde}および各\texttt{.ino}を使用した。
  \item ドラムは成毛がGitHubへアップロードしたファイルを輪唱用に改造して組み込んだ（佐藤が後ほどGitHubへアップロード予定）。
  \item 輪唱が成功した場合，以下のタスクは完了とみなすこととした。
  \begin{itemize}
    \item 2.3 共通ライブラリ整備（PROGMEM＋I2Cフレーム）
    \item 2.2.2 CONFIG送信処理（起動時 entry\_offset配布）
    \item 2.3.1 ISR受信・デコード処理（onReceive）
  \end{itemize}
  \item 「1.3 筐体制作・実装」と「2.2.1 BPM変化制限ロジック（1小節最大10BPM）」には依存関係があり，前者が後者の前提となるため，早急に1.3の筐体制作へ着手する方針とした。
  \item 本日，以下の5項目のテストを実施する（前提：輪唱可能かつ楽器音OK，PLL／BPM対応は未実装）。
  \begin{itemize}
    \item 楽器音の同期誤差
    \item パケットロス率
    \item 楽曲識別
    \item 楽器識別（4種類）
    \item リラクゼーション感のチェック
  \end{itemize}
\end{itemize}

\subsubsection*{議論内容}
\begin{itemize}[leftmargin=1.5em]
  \item 今週の成果として，各楽器音の精度向上および通信処理の改善を共有した。
  \item ファイル構成について山口より共有があった。各楽器のスレーブ用\texttt{.ino}は楽器ごとのフォルダ内\texttt{arduino}に，各楽器の\texttt{.pde}は同フォルダ内\texttt{Processing}に格納している。マスター用ファイルは\texttt{sync}フォルダ内の\texttt{arduino}に格納されている。
  \item PLL的補正は，可変テンポ対応と並行して実装する方針である。
  \item 6月3日は集まることができず実機による同期テストを実施できなかったため，来週までに実施する。中間テスト等が重なりメンバー間の連絡が滞っていたため，今週からは活発にコンタクトを取る。
  \item ハード設計（筐体）が輪唱の開始トリガを兼ねており，この遅延が依存タスク全体を停滞させている。それ以外はスケジュールどおりであるため，今週のタスクはハード設計を中心とする。
\end{itemize}

\subsubsection*{ToDo / アクション}
\begin{itemize}[leftmargin=1.5em]
  \item 改造済みドラムファイルのGitHubアップロード……佐藤
  \item 1.3 筐体制作・実装（最優先・後続タスクの依存元）……細澤
  \item 2.2.1 BPM変化制限ロジック（1小節最大10BPM，1.3完了後）
  \item Serial制御ロジック実装（5台同時受信・パート振分け）
  \item 2.2.5 サーボ制御・状態管理実装（IDLE→PLAYING→FINISHED）
  \item 3.1 単体・同期精度テスト（V1聴取＋V2 FFT比較）
  \item 1.5 ロードセルの反応検証
  \item Slack上の会話履歴のZulipへの移行……酒井・細澤
\end{itemize}

\subsubsection*{保留・次回検討事項}
\begin{itemize}[leftmargin=1.5em]
  \item 第6週目までの進捗を踏まえ，最終成果物の最低限ライン（MVP）を更新する。
  \item ポテンショメータ関連のタスクが未着手であり，着手時期を検討する。
\end{itemize}

%-------------------------------------------------------------------
\subsection*{議題②：計画からの変更点}

\subsubsection*{決定事項}
\begin{itemize}[leftmargin=1.5em]
  \item 計画書・スケジュールそのものからの変更は現状なし。
  \item ロードセルによる開始処理が未実装のため，現状はボタンで開始・終了を行う運用で代替する。
\end{itemize}

\subsubsection*{議論内容}
\begin{itemize}[leftmargin=1.5em]
  \item ハード設計（筐体）が開始トリガを兼ねており，その遅延が依存タスクを含めた遅れを生んでいる。当面はボタン操作で代替し，最低保証ラインを確保する。
\end{itemize}

\subsubsection*{ToDo / アクション}
特になし。

\subsubsection*{保留・次回検討事項}
\begin{itemize}[leftmargin=1.5em]
  \item 最低保証ライン：ボタンで開始・終了を行う輪唱システムは提供可能。
\end{itemize}

%-------------------------------------------------------------------
\subsection*{議題③：個人毎の次回までに実施する内容}

\subsubsection*{決定事項}
各メンバーが次回（6月17日）までに実施する内容を以下に示す。
\begin{itemize}[leftmargin=1.5em]
  \item 山口……各楽器の精度向上，PLL的補正・可変テンポ対応の実装
  \item 佐藤……改造済みドラムファイルのGitHubアップロード，同期精度テストの実施
  \item 成毛……オルガン・パーカッションの作成継続（パーカッションは2種のうち1種が完成済み）
  \item 細澤……1.3 筐体制作・実装，2.2.1 BPMロジックの完成，Slack$\to$Zulip移行, Github進捗管理, 
  \item 武本……1.3 筐体制作・実装，2.2.1 BPMロジックの完成
  \item 酒井……1.3 筐体制作・実装，2.2.1 BPMロジックの完成, 本日実施した5項目テスト結果の取りまとめ，Slack$\to$Zulip移行
\end{itemize}

\subsubsection*{議論内容}
\begin{itemize}[leftmargin=1.5em]
  \item 次回までの最優先事項は，遅延タスク（筐体制作＋BPMロジック）の完了，およびロードセルを検知してベビーメリーの回転までを実現すること，加えて5台同時のシリアル制御を確立することである。
\end{itemize}

\subsubsection*{ToDo / アクション}
議題①のとおり。

\subsubsection*{保留・次回検討事項}
特になし。

%-------------------------------------------------------------------
\section*{4.次回予定}
\begin{tabular}{ll}
  \toprule
  日時・場所 & 2026年6月17日（水）・731教室 \\
  議題       & 遅延タスク（筐体制作・BPMロジック）の完了報告，5台同時シリアル制御，実機同期テスト結果 \\
  その他，特記事項 & 今週から活発にコンタクトを取る \\
  \bottomrule
\end{tabular}

\end{document}
